\documentclass[../Main.tex]{subfiles}


\begin{document}

\section{Theoretischer Hintergrund \& Analyseframework}

\subsection{Definition des Hack-and-Slash Genre}

In der Videospielindustrie bezeichnet "Hack and Slash" eine Subgenre der Action Games. Der Fokus liegt dabei auf Echtzeitkampf mit Nahkampfwaffen, wie Schwertern, Äxten oder anderen Klingen. Fernkampfwaffen sind zwar vertreten, spielen aber nur eine untergeordnete und unterstützende Rolle. Das Genre entstammt den 2D Beat'em Up Games, entwickelte sich weiter in der 3D Ära und wurde vor allem durch Titel wie Devil May Cry (2001) oder Bayonetta (2009) geprägt.

Die modernen Spiele zeichnen sich durch hohe Spielgeschwindigkeiten und stilvolle, fliessende Bewegungen aus. Das Verketten präziser Angriffs-Kombos, das gezielte Abbrechen von Bewegungen, I-Frames und strategisches Distanzmanagement zum Gegner zeichnen das Gameplay aus. Das, was ein überzeugendes Hack-and-Slash ausmacht, ist nicht allein die Anzahl dieser Mechaniken, sondern das Zusammenspiel dieser. Erst wenn präzises Kampfverhalten und befriedigendes Spielgefühl aufeinander abgestimmt sind, entsteht das genredefinierende Gefühl von Kontrolle und Ausdrucksstärke.

Um dieses Zusammenspiel systematisch zu erfassen und vergleichbar zu machen, bedarf es eines geeigneten Analyserahmens. Im folgenden Abschnitt wird das MDA-Framework als methodisches Werkzeug eingeführt, das die Beziehung zwischen Spielregeln, emergenten Verhaltensmustern und dem emotionalen Spielerleben strukturiert beschreibt.


\subsection{Das MDA-Analyseframework}

Das MDA-Framework (Mechanics, Dynamics, Aesthetics) wurde von Hunicke, LeBlanc und Zubek entwickelt und ist ein formaler Ansatz zur Analyse und zum Verständnis von Spieldesign. Es beschreibt Spiele nicht als statisches Medium wie Buch oder Film, sondern als interaktives Erlebnis, welches aus dem Zusammenspiel dreier klar abgegrenzter Ebenen entsteht.

Die unterste Ebene, die Mechaniken, umfasst die konkreten Regeln und Systeme eines Spiels. Darunter versteht man alles, was durch Daten und Algorithmen getrieben ist, wie Schadenswerte, Treffererkennung, Bewegungsgeschwindigkeit oder die Gegner-KI. Daraus entstehen zur Laufzeit die Dynamiken, also das emergente Verhalten, das sich aus dem Zusammenspiel von Mechaniken und Spielerentscheidungen ergibt. Konkret geht es um den Kampfrhythmus, taktische Ausweichmuster oder eskalierende Gegnerwellen. Die oberste Ebene, die Ästhetik, beschreibt schliesslich die emotionale Reaktion des Spielers auf dieses System. Gefühle wie Kontrolle und Befriedigung bei gelungenen Combos stehen im Vordergrund.

Entscheidend ist die Kausalität zwischen den Ebenen. Mechaniken erzeugen Dynamiken, die wiederum Ästhetik erzeugen. Gleichzeitig erleben die Spieler das Game in umgekehrter Reihenfolge. Für die Analyse bestehender Spiele bedeutet das, dass man vom Spielgefühl beginnt und sich systematisch zu den Ursachen, also den Mechaniken, vorarbeiten muss.

Dieses Framework eignet sich besonders gut für die Analyse von Hack-and-Slash Spielen, da das Genre ein besonders hohes Mass dieser Kausalität aufweist. Beispielsweise sind die genauen Parameter einer einzelnen Mechanik, wie einer Attacke, bedeutungslos, wenn sie nicht im Kontext weiterer Systeme, wie Gegnerlebenspunkten, Gegnerverhalten und Animationsgeschwindigkeit, betrachtet werden.\autocite{hunicke2004mda}

Auf Basis dieses Frameworks wurde ein strukturierter Fragebogen mit 15 Fragen entwickelt, der die drei Ebenen getrennt betrachtet. Er dient als Grundlage für die Analyse der Referenzspiele ab den Abschnitt \ref{sec:Referenzspielanalyse} sowie als inhaltliche Basis für den Playtest-Fragebogen.


\subsection{Analysefragebogen}

Die 15 Fragen sind den drei MDA-Ebenen zugeordnet und werden in umgekehrter Reihenfolge gestellt, da die Analyse beim Spielgefühl beginnt und sich systematisch zu den Ursachen vorarbeitet. Thematisch sind die Fragen so ausgerichtet, dass sie gezielt Gegnerverhalten, Gegnerplatzierung, Spielfluss, Schwierigkeit und Feedback beleuchten, da diese Aspekte im Zentrum unserer Erkenntnisfrage stehen.

\subsubsection{Aesthetics (Ästhetik und emotionales Erleben)}

\begin{enumerate}[label=\textbf{A\arabic*}, leftmargin=1.5cm, itemsep=0.5em]
    \item Welche emotionale Reaktion löst der Kampf primär aus: Kontrolle, Befriedigung, Anspannung oder Frustration? Warum?
    \item Wie gut vermittelt das audiovisuelle Feedback die Wucht und Präzision der Treffer? Wie?
    \item Vermittelt die Gegnerplatzierung ein Gefühl von intentionalem Design, wirkt die Anordnung der Gegner durchdacht oder zufällig und wie beeinflusst diese das Empfinden des Spielers?
    \item Verändert sich das emotionale Erleben im Verlauf des Levelabschnitts und ist eine Spannungskurve erkennbar?
    \item Unterscheiden sich die Combos darin, welche Emotionen sie auslösen, bringen sie den Spieler in Bedrängnis oder gibt es ein Gefühl der Erleichterung?
\end{enumerate}

\subsubsection{Dynamics (Spieldynamik und Interaktion)}

\begin{enumerate}[label=\textbf{D\arabic*}, leftmargin=1.5cm, itemsep=0.5em]
    \item Wie sieht der Levelabschnitt von Beginn bis Ende aus, erläutere erkennbare Muster und wie etwaige mit der Spannungskurve zusammenhängen?
    \item Erzeugt die Kombination verschiedener Gegnertypen innerhalb eines Encounters taktischen Druck oder agieren die Gegner unabhängig voneinander?
    \item Entsteht durch Gegnerplatzierung/Spawnpunkte und deren Anzahl ein flüssiger Spielfluss, gibt es einen spürbaren Unterschied zwischen Einzel- und Gruppenkonfrontationen?
    \item Wie erholt sich der Spieler nach einem Fehler, ist die Feedback-Schleife zwischen Fehler und Korrektur nachvollziehbar?
    \item Gibt es verschiedene Combos/Kampfabläufe und folgen daraus verschiedene Arten die Gegner zu bekämpfen?
\end{enumerate}

\subsubsection{Mechanics (Spielmechaniken)}

\begin{enumerate}[label=\textbf{M\arabic*}, leftmargin=1.5cm, itemsep=0.5em]
    \item Wie unterscheiden sich die Gegnertypen auf Mechanikenebene: Bewegungsgeschwindigkeit, Angriffsreichweite, Schadenswerte?
    \item Wie ist das Angriffsverhalten der Gegner implementiert, reagieren sie aktiv auf den Spieler oder folgen sie festen vordefinierten Mustern?
    \item Wie begrenzt das Leveldesign die Bewegungsoptionen von Spieler und Gegner, gibt es räumliche Einschränkungen die den Kampf beeinflussen?
    \item Wie ist das Hitstun-System und Knockback-System designed, lassen sich verschiedene Arten beobachten?
    \item Wie hängen die gedrückten Tasten mit dem visuellen Kampfereignissen zusammen, gibt es eine bestimmte Intention hinter den Tasten?
\end{enumerate}


\subsection{Analyse: Ninja Gaiden Sigma Plus} \label{sec:Referenzspielanalyse}

Im folgenden Abschnitt werden Kampfsequenzen aus dem Spiel \textit{Ninja Gaiden Sigma Plus} mit Hilfe des MDA-Fragebogens analysiert. Das Spiel wurde von \textit{Team Ninja}, einem Studio des Japanischen Videospielentwickler \textit{Koei Tecmo}, im Jahre 2012 für die \textit{PlayStation Vita} entwickelt.\autocite{wikipedia_ninjagaiden_2026} Dafür haben wir die \textit{Ninja Missions}, eine Kollektion von längeren Combat Sequenzen gespielt, betrachtet und obriges Analyse-Framework verwendet.

\subsubsection{Inhaltlicher Kontext}

In \textit{Ninja Gaiden Sigma Plus} steuert man dne Protagonisten in der Third-Person-Perspektive. Dieser ist der Ninja \textit{Ryu Hayabusa}. Wie auch in anderen Action-RPGs funktioniert die Metaprogression des Spiels durch die Erweiterung der Fähigkeiten des Spielers. Diese findet durch das Fortschreiten des Spielgeschehens und dem daraus folgenden Erlangen von Gegenständen statt. Dazu gehören Items, welche die Lebenspunkte des Spielers permanent steigern können, Artefakte, die den Spieler in einem bestimmten Aspekt stärken, auffindbare Waffen, Erwerbbare Waffen-Upgrades oder sogenannte \textit{Technique Scrolls}, welche eine Combo freischalten. Ein Level wird bewältigt indem man ein deffiniertes Ziel erreicht, z.B. durch das Erreichen eines bestimmten Punktes im Level oder einen herausfordernder Kampf gegen einen \textit{Bossgegner}. Während dem Level treten diverse Hindernisse auf, unteranderem physische Barrieren, die durch eine Schlüssel-Schloss-Interaktion beseitigt werden können oder Kampfbegegnungen mit Feinden. Ryu kann seine Gegner mit verschiedenen Waffen und deren Combos angreifen. Die bekannteste Combo in der Ninja Gaiden Franchise ist die \textit{Izuna Drop}, wobei der Protagonist den Feind in die Luft wirft und Kopfüber auf den Boden schlägt.


\subsubsection{Asthetik}

\begin{description}
    \item[\textbf{A1}] \textbf{Welche emotionale Reaktion löst der Kampf primär aus: Kontrolle, Befriedigung, Anspannung oder Frustration? Warum?}
\end{description}

Ninja Gaiden gibt dem Spieler grosse Kontrolle über alle Situationen, kann diese jederzeit aber auch wegnehmen. Bewegungen und Knopfdrücke sind sofortig und präzise, und das Spiel gibt dem Spieler jedes Werkzeug für jede Situation. Wenn man von zu vielen Gegnern umgeben ist, kann man eine schnelle Rolle ausführen. Der Spieler kann mit einem rapiden Sturzangriff sofort auf weiter entfernte Gegner einstechen, mit schweren Angriffen die Verteidigung der Gegner durchbrechen und mit besonders vielen Angriffen den Gegner zu Boden stürzen lassen, wo dieser hilflos liegt, bis er wieder aufsteht. Der Spieler besitzt nicht nur Kontrolle auf dem Boden, sondern auch in der Luft. Hat man einen Gegner in die Luft geschlagen, so ist dieser hilflos. Dies ist befriedigend, weil der Spieler in diesem Moment absolute Macht über den Feind besitzt und selbst entscheiden kann, wie er die Situation auflöst. Dies ist auch der Grund, wieso der Izuna Drop eine Ikonen-Combo unter Action Games ist: Sie ist extrem befriedigend auszuführen, weil sie dem Spieler ein Gefühl absoluter Kontrolle und Macht über den Feind vermittelt.

Doch wenn man zu rücksichtslos agiert und einen Treffer kassiert, nimmt das Spiel sofort die Kontrolle weg. Erkennt man nicht rechtzeitig, dass man von Gegnern umzingelt ist, schlagen diese sofort zu. Sobald ein Gegner beginnt, den Spieler mit einer eigenen Combo anzugreifen, lässt das Spiel den Spieler erst entkommen, wenn der Gegner fertig ist. Wird der Spieler von einer besonders schweren Attacke getroffen, kann er nach hinten geschleudert werden und muss erst wieder aufstehen, bevor er die Kontrolle zurückgewinnt. Dies ist eine der Weisen, wie das Spiel den Spieler für übermässige Aggression bestraft. Doch das Spiel bestraft auch einen übermässig passiven Spielstil. Die meisten Gegner können die Verteidigung des Spielers durchbrechen, wenn sie ihn oft genug angreifen. Ausserdem besitzen manche Gegner einen Angriff, der den Spieler in eine eindrucksvolle Animation gefangen hält, bis diese vorbei ist. Diese Angriffe sind spektakulär anzusehen, doch der Gegner ist derjenige, der das Spektakel betreibt und den Hauptcharakter dabei demütigt.

Diese Momente fühlen sich zwar frustrierend an, jedoch nicht unfair. Wie bereits erläutert gibt das Spiel dem Spieler jedes Werkzeug für jede Situation, und man erkennt, dass es die eigene Schuld ist und man zu weit gegangen ist. Genau diese Frustration lässt Momente der Kontrolle sich noch besser anfühlen. Wenn der Spieler wieder aufsteht, seinen Fehler erkennt, sich verbessert und daraufhin wieder die Hand über die Situation gewinnt, bringt dies immense Befriedigung. Diese dünne Linie zwischen absoluter Kontrolle und totalem Kontrollverlust erzeugt die Spannung, die den Kampf in Ninja Gaiden so packend macht.

\begin{description}
  \item[\textbf{A2}] \textbf{Wie gut vermittelt das audiovisuelle Feedback die Wucht und Präzision der Treffer? Wie?}
\end{description}

Es hat keinen Moment gegeben, in dem man das Gefühl hatte, der Angriff hätte treffen sollen, dies aber nicht passiert ist. Jeder Treffer ist gelandet und hat sich angemessen angefühlt. Das ist befriedigend, weil das audiovisuelle Feedback den Spieler unmittelbar bestätigt, dass seine Eingabe korrekt war. Ein erfolgreicher Treffer wird stets durch ein helles, spritzendes Geräusch gekennzeichnet. Der Gegner begibt sich sofort in eine rückwärtslehnende Position, Funken und Blut spritzen. Der Aufprallklang passt dabei zur jeweiligen Waffe: Ein Schwert klingt wie das Reissen von Metall durch Fleisch, ein Nunchaku wie ein harter Aufprall. Schlägt man hingegen auf einen blockierenden Feind ein, wird ein dumpfes Geräusch erzeugt, und anstelle von Funken erscheint eine kleine Schockwelle. Bricht man die Blockade, stolpert der Gegner nach hinten, die Schockwelle zersplittert und ein besonders klares Geräusch von Metall auf Metall erklingt. Dies ist befriedigend, weil der Spieler durch das veränderte Feedback unmittelbar erkennt, dass er eine Schwelle überwunden hat.

Der einzige, jedoch gewichtige Kritikpunkt betrifft den Izuna Drop. Der abschliessende Aufprall auf den Boden sollte, wie bei den meisten anderen Angriffen im Spiel, die einen Feind aus der Luft zu Boden schleudern, in einem Blutbad enden. Stattdessen erklingt lediglich ein dumpfer Knall, was angesichts der ikonischen Bedeutung dieser Combo enttäuschend wirkt.

\begin{description}
  \item[\textbf{A3}] \textbf{Vermittelt die Gegnerplatzierung ein Gefühl von intentionalem Design, wirkt die Anordnung der Gegner durchdacht oder zufällig und wie beeinflusst diese das Empfinden des Spielers?}
\end{description}

Die Gegnerplatzierung in den Ninja Missions ist grösstenteils zufällig generiert, was das Spielerlebnis jedoch kaum beeinträchtigt. Die Ninja Missions spielen sich vorwiegend in engen Räumen ab, in denen eine feste Platzierung ohnehin wenig Mehrwert böte. Das Spiel platziert neue Feinde stets ausserhalb des Sichtfelds der Kamera oder an Stellen, die visuell glaubwürdig wirken, etwa durch Löcher im Dach, wie in Ninja Mission 1.1. Diese Regel trägt dazu bei, die offensichtliche Spielhaftigkeit zu verbergen und den Spieler in der Spielwelt zu halten, was sich positiv auf die Immersion auswirkt.

Das Fehlen fester Spawnpunkte schadet dem Spielerlebnis auch deshalb kaum, weil die Gegner-KI von Ninja Gaiden äusserst dynamisch ist. Gegner mit Fernkampfwaffen halten gezielt Distanz zum Spieler, während aggressivere Feinde selten stillstehen und sich genauso schnell durch das Level bewegen wie der Spieler selbst. Es kommt häufig vor, dass ein Feind hinter den Spieler sprintet, sobald dieser sich zu stark auf einen anderen Gegner konzentriert. Dieser Moment fühlt sich erschreckend und gleichzeitig fair an, weil er auf eine eigene Unachtsamkeit zurückzuführen ist. Zudem greifen nahestehende Gegner oft gemeinsam an, sobald die Verteidigung des Spielers gebrochen ist, was verheerende Konsequenzen haben kann. All dies wirkt keineswegs zufällig, sondern äusserst intentionell. Die dynamische Organisation der Gegner erzeugt ein Gefühl von taktischem Druck, das den Kampf spannend hält.

\begin{description}
  \item[\textbf{A4}] \textbf{Verändert sich das emotionale Erleben im Verlauf des Levelabschnitts und ist eine Spannungskurve erkennbar?}
\end{description}

Die Spannung in einem Kampf ist nicht konstant, sondern verändert sich in Abhängigkeit von der Anzahl der aktiven Gegner, ihrer individuellen Schwierigkeit sowie ihrer relativen Schwierigkeit im Vergleich zur vorherigen Welle. Generell lässt sich sagen, dass Kämpfe ihre Spannung hauptsächlich durch eine schrittweise Steigerung von Gegneranzahl und Gegnerschwierigkeit aufbauen. Dies ist befriedigend, weil der Spieler stets das Gefühl hat, dass der Druck kontrolliert eskaliert und nicht willkürlich zunimmt. Er kann sich auf die Kurve einstellen und wird dafür belohnt, wenn er sie meistert.

\begin{description}
  \item[\textbf{A5}] \textbf{Unterscheiden sich die Combos darin, welche Emotionen sie auslösen, bringen sie den Spieler in Bedrängnis oder gibt es ein Gefühl der Erleichterung?}
\end{description}

Das \textit{Dragon Sword}, eine der Hauptwaffen im Spiel, besitzt rund 10 Combos sowie mehrere eigenständige Aktionen. Die meisten davon lösen entweder keine oder nur schwache Emotionen aus; einige wenige hingegen erzeugen grosse Befriedigung. Die Combos des Dragon Sword folgen alle dem Muster zuerst leichter und danach Schwerer Angriff:

\begin{equation*}
n * \text{Leichter Angriff} + m * \text{Schwerer Angriff}
\end{equation*}

Dies führt zwar zu Combos, die leicht zu merken sind, jedoch in ihrer Differenzierung schwach ausfallen. Selbst wenn die Animationen stark variieren, hat man oft das Gefühl, ungefähr dieselben Knöpfe gedrückt zu haben. Gelegentlich passiert es zudem, dass die ausgeführte Combo nicht die gewünschte war, was zwar das Spielerlebnis nicht stark beeinträchtigt, aber ein Gefühl der Ungenauigkeit hinterlässt. Die mechanischen Hintergründe dieses Gefühls werden später weiter ausgeführt.

Eigenständige Aktionen entstehen grundsätzlich aus der Kombination eines Angriffs mit einem Richtungs-Input. Zwei Aktionen stechen dabei besonders heraus. Die erste ist der \textit{Flying Swallow}: ein rasanter Sturzflug aus der Luft, bei dem der Protagonist auf den nächsten Feind zufliegt und diesen mit einem Hieb, oft einer Enthauptung, trifft. Diese Aktion ist befriedigend, weil sie dem Spieler das Gefühl von Geschwindigkeit und Präzision vermittelt und gleichzeitig grossen Schaden anrichtet. Die zweite Aktion ist kontextabhängig und besteht aus mehreren Teilbewegungen: Durch langes Gedrückthalten der Taste für einen schweren Angriff wird ein Spezialangriff aufgeladen. Jede Waffe besitzt hierfür eine eigene Animation; beim Dragon Sword entfaltet sich eine Sequenz mehrerer harter Schnitte auf den Gegner, während der Spieler weder Schaden nehmen noch der Gegner sich verteidigen kann. Die Kamera rückt dabei nah heran. Dies ist befriedigend, weil der Spieler für einen Moment vollständige Kontrolle und Unverwundbarkeit erlebt, ein bewusstes Belohnungsmoment des Designs.

Zusammenfassend lässt sich sagen, dass viele Combos und Aktionen in \textit{Ninja Gaiden Sigma Plus} für den durchschnittlichen Spieler nicht zwingend notwendig für den Spielfortschritt sind. Sie bringen den Spieler in der Regel nicht in Bedrängnis, und wer weiss, wie man sie korrekt einsetzt, kann das Spielerlebnis deutlich verbessern. Die Umsetzung der Tastenkombinationen bleibt dennoch leicht unbefriedigend. Während das Dragon Sword rund 10 Combos besitzt, schreibt das Spiel der Waffe \textit{Dragon's Claw and Tiger's Fang} rund 27 zu, eine Komplexität, die den Spieler durchaus in Bedrängnis bringen kann, da die effiziente Nutzung dieser Waffe ohne Kenntnis ihrer Mechaniken erheblich erschwert wird. Diese zu erlernen erfordert einen grossen Zeitaufwand, was für den durchschnittlichen Spieler eine echte Hürde darstellt.


\subsubsection{Dynamik}

\begin{description}
  \item[\textbf{D1}] \textbf{Wie sieht der Levelabschnitt von Beginn bis Ende aus, erläutere erkennbare Muster und wie etwaige mit der Spannungskurve zusammenhängen?}
\end{description}

Eine Kampfsequenz in Ninja Gaiden Sigma wird in einer Wellenstruktur aufgebaut. Eine Kampfsequenz kann mehrere Wellen definieren, indem eine bestimmte Anzahl einer bestimmten Art von Gegner auftaucht. Zudem kann eine Welle auch einen Grenzwert von Gegnern definieren, die gleichzeitig auf dem Feld sind. Dieser liegt zumeist bei maximal vier Gegnern. Sobald ein Feind gestorben ist, wird dieser sofort ersetzt, solange in der Definition der Welle noch Gegner dieser Art vorhanden sind.

Diese Struktur bildet eine Spannungskurve, die beim anfänglichen Auftauchen der Gegner und deren darauffolgender Ersetzung konstant steigt, aber langsam sinkt, sobald die Anzahl der Gegner auf dem Feld abnimmt. Dies bedeutet, dass sich der Spieler am Ende einer Welle leicht erholen kann, bevor die nächste beginnt. Ausserdem werden die Wellen während einer Kampfsequenz tendenziell schwieriger, entweder durch den Einsatz von Gegnern mit komplexerer KI und mehr Leben oder durch die Steigerung der Gegneranzahl. Ohne diese Wellenmechanik würde die Spannungskurve bei längeren Kämpfen extrem rasch steigen, und der Spieler wäre schnell mental überlastet, was zu mehr Fehlern und schlechteren Entscheidungen führen würde.

\begin{description}
  \item[\textbf{D2}] \textbf{Erzeugt die Kombination verschiedener Gegnertypen innerhalb eines Encounters taktischen Druck oder agieren die Gegner unabhängig voneinander?}
\end{description}

Wie bereits erläutert, agieren Gegner durchaus nicht unabhängig voneinander, und sollte dies eine Illusion sein, ist diese äusserst überzeugend. Generell können Gegner in zwei Kategorien eingeteilt werden: Fern- und Nahkämpfer. Das Verhalten der Nahkämpfer ist äusserst divers und wird ebenfalls in der Sektion über Mechaniken erläutert. Manche Gegner sind zudem in der Lage, sowohl aus der Ferne als auch aus der Nähe anzugreifen. Dieses Verhalten wird eingesetzt, sobald der Gegner sich etwas zu weit vom Spieler entfernt befindet.

Die Kombination von Fern- und Nahkampfgegnern in einem Kampf erzeugt taktischen Druck, indem sie den Spieler zu einer Entscheidung zwingt. Man kann sich dazu entscheiden, zuerst den Fernkampfgegner zu beseitigen, riskiert dabei aber, von oft stärkeren Nahkampfgegnern bedrängt zu werden und einen verheerenden Rückentreffer zu kassieren. Entscheidet sich der Spieler dagegen und geht zuerst gegen die Nahkampfgegner vor, begeht er ebenfalls ein Risiko. Ein zeitlich ungünstiger Angriff aus der Ferne, während der Spieler sich im aktiven Kampf befindet, kann diesen mitten aus einer Combo herausreissen und direkt in eine Barrage von Gegnerangriffen werfen. Dies ist befriedigend, weil beide Entscheidungen mit echten Konsequenzen verbunden sind und der Spieler das Ergebnis seiner Wahl unmittelbar spürt.

\begin{description}
  \item[\textbf{D3}] \textbf{Entsteht durch Gegnerplatzierung/Spawnpunkte und deren Anzahl ein flüssiger Spielfluss, gibt es einen spürbaren Unterschied zwischen Einzel- und Gruppenkonfrontationen?}
\end{description}

In Ninja Gaiden Sigma ist der Kampf gegen einen einzelnen Gegner im Vergleich zu einem Gruppenkampf trivial. Der Spieler hat alle Zeit der Welt, um auf den Gegner zu reagieren, sich zu verteidigen und im richtigen Augenblick zurückzuschlagen. Der Tanz zwischen Spieler und Feind ist langsam, selbst wenn das Kampfsystem, auf dem er aufbaut, methodisch und durchdacht ist. Doch wie bereits erläutert kann selbst die Einführung eines einzigen weiteren Feindes zu einem drastischen Anstieg der Schwierigkeit führen. Es braucht dabei nicht immer eine Kombination aus Nah- und Fernkampfgegner, um diesen Effekt zu erzeugen. Auch der Kampf gegen mehrere Nahkampfgegner ruft ein ähnliches Entscheidungsspiel hervor. Schwache Gegner können den Spieler von der Seite bedrängen und überraschen, während ein stärkerer Gegner dem Spieler ebenfalls keine Luft zum Atmen lässt. Wie vieles in diesem Spiel erfordert dies blitzschnelle Entscheidungen und eine Balance von Risiko und Belohnung. Mehrere gleichzeitig agierende Gegner haben somit den Effekt einer konsequenten Beschleunigung des Spielflusses.

Die bereits erläuterte Umzingelung des Spielers durch mehrere Gegner trägt ebenfalls zu einem schnelleren und flüssigeren Spielfluss bei. Würden alle Gegner direkt vor dem Spieler stehen, könnte dieser jederzeit nach hinten ausweichen. Doch dieses Gegnerverhalten verhindert diesen einfachen Ausweg. Der Spieler muss ständig wachsam sein, nicht nur gegenüber den Feinden vor ihm, sondern auch hinter ihm und gegenüber allen Gegnern, die ihn potentiell einkreisen könnten. Er wird sich seiner Umgebung bewusst und muss jederzeit bereit sein, sich aus einer potentiellen Zwickmühle herauszurollen. Diese Mechanik erzeugt enorme mentale Belastung, und es ist äusserst befriedigend, wenn man sich durch das Netz von Gegnern hindurchfechtet. Nach längerer Zeit wird sie jedoch anstrengend, was ein weiterer Grund für die Notwendigkeit der Wellenmechanik ist.

\begin{description}
  \item[\textbf{D4}] \textbf{Wie erholt sich der Spieler nach einem Fehler, ist die Feedback-Schleife zwischen Fehler und Korrektur nachvollziehbar?}
\end{description}

Ninja Gaiden Sigma unterscheidet im Kampf zwischen zwei Schadenszuständen. Der erste ist ein kurzes Zucken nach einem Treffer, der zweite ist ein Sturz auf den Boden. Sowohl Spieler als auch Gegner besitzen diese zwei Zustände, und die zugrunde liegende Mechanik funktioniert für beide ähnlich. Nach einem kurzen Treffer kann man sich relativ schnell erholen, während ein Sturz auf den Boden eine längere Aufsteh-Animation erfordert. Wird man in der Luft getroffen, wechselt man direkt in den Sturzzustand, was dazu dient, die Stärke von Luftbewegungen auszubalancieren. Bei kleinen Fehlern erholt sich der Spieler schnell, bei grösseren Fehlern, bei denen man sich mehrmals treffen lässt, dauert die Erholung etwas länger. Die Bestrafung fühlt sich jedoch fair an, da jeder Treffer auf den Spieler dieselben visuellen Effekte erzeugt wie ein Treffer auf einen Feind. Zudem ist die Erholung automatisch und zeitlich angemessen, da der Spieler nicht lange in der Animation gefangen bleibt. Gegner greifen den Spieler nach einem Sturz ausserdem nicht exzessiv an. Selbst wenn das Spiel die Leistung des Spielers streng bewertet, weiss es auch Grenzen zu ziehen.

\begin{description}
  \item[\textbf{D5}] \textbf{Gibt es verschiedene Combos/Kampfabläufe und folgen daraus verschiedene Arten die Gegner zu bekämpfen?}
\end{description}

Die verschiedenen Combos und Aktionen, erläutert in Frage A5, sind die Lösungen zu einem Puzzle, wobei das Puzzle den Kampf darstellt. Darunter lassen sich mehrere generelle Muster erkennen, die das Spiel regelmässig einsetzt.

Combos, die mit einem schweren Angriff beendet werden, können die Verteidigung eines Gegners durchbrechen, während dies mit mehreren aufeinanderfolgenden leichten Angriffen allein nicht möglich ist. Beide Kampfweisen sind viable, jedoch werden erfahrenere Spieler die erste Option stark bevorzugen. Einen Verteidigungs-durchbrechenden Angriff gibt es nicht nur am Boden, sondern auch in der Luft. Der \textit{Helmet Breaker} ist eine Aktion, bei der der Protagonist aus der Luft mit grosser Wucht auf den Boden stürzt. Das Spiel gibt Spielern damit sowohl in der Luft als auch am Boden die gleichen Möglichkeiten, und man kann aus beiden Positionen heraus den Gegner bedacht besiegen.

Der Flying Swallow, die schnellen Sprünge, die gute Luftmobilität und die raschen Ausweichrollen unterstützen einen beweglichen Spielstil. Die aufladbaren Spezialangriffe sowie schwere und explizite Combos schaffen im Gegensatz dazu einen methodischen, bodenbasierten Spielstil. Der Izuna Drop ist die Combo, in der Luft und Boden, das Agile und das Methodische zusammenkommen. Der rasche Wurf des Gegners in die Luft repräsentiert das Agile, während die notwendige Vorbereitung, um zuvor die Verteidigungshaltung des Gegners zu durchbrechen, das Methodische repräsentiert.

Zudem stellt der aufladbare Spezialangriff einen weiteren Kampfstil zur Verfügung. Die volle Aufladung dauert rund fünf Sekunden, während derer sich der Spieler weder bewegen noch verteidigen kann, ausser er bricht den Angriff ab. Fünf Sekunden sind eine erhebliche Zeitspanne für ein Spiel, das eigentlich grossen Wert auf sofortige Reaktionen legt. Dieser Angriff stellt daher einen High-Risk-High-Reward-Spielstil dar, den das Spiel für jede Waffe anbietet. Zudem kann der Angriff sofort ausgeführt werden, wenn in der Umgebung \textit{Essence} vorhanden ist. Essence sind kleine leuchtende Kugeln, die je nach Farbe entweder als Währung oder Heilmittel dienen. Die Essence wird dabei verbraucht, und sollte ein Feind den Spieler treffen, bevor dieser den aufgeladenen Angriff einsetzen kann, hat der Spieler nicht nur Schaden genommen, sondern auch die Essence dauerhaft verloren. Der Spieler muss also auch in dieser Hinsicht eine Entscheidung treffen, die nicht direkt mit den unmittelbaren Kampfereignissen zusammenhängt: ob er die Essence einfach normal aufnehmen oder riskieren soll, den Spezialangriff aufzuladen.


\subsubsection{Mechanik}

\begin{description}
  \item[\textbf{M1}] \textbf{Wie unterscheiden sich die Gegnertypen auf Mechanikenebene: Bewegungsgeschwindigkeit, Angriffsreichweite, Schadenswerte?}
\end{description}



\begin{description}
  \item[\textbf{M2}] \textbf{Wie ist das Angriffsverhalten der Gegner implementiert, reagieren sie aktiv auf den Spieler oder folgen sie festen vordefinierten Mustern?}
\end{description}



\begin{description}
  \item[\textbf{M3}] \textbf{Wie begrenzt das Leveldesign die Bewegungsoptionen von Spieler und Gegner, gibt es räumliche Einschränkungen die den Kampf beeinflussen?}
\end{description}



\begin{description}
  \item[\textbf{M4}] \textbf{Wie ist das Hitstun-System und Knockback-System designed, lassen sich verschiedene Arten beobachten?}
\end{description}

Hinter den bereits erläuterteten Trefferarten, nämlich des Zucken und dem Sturz, steckt warscheinlich ein versteckter Wert. Dieser wird oft als \textit{Super-Armor} bezeichnet. Im Gegensatz zu einer konventionalen Rüstung in einen Spiel beschützt die Super-Armor den Spieler nicht direkt, denn die Super-Armor ist funktionell getrennt von der Lebensleiste. Jeden Angriff kann eine gewisse Anzahl von Lebenspunkten wegnehmen, aber auch einige Super-Armor Punkte. Solange diese nicht leer ist, wird der Charakter nur leicht Zucken, sobald sie leer ist, stürzt man auf den Boden. Spieler und Feinde besitzen beide diese Mechanik. Bei dem Rückstoss können Waffen zusätzliche auch eine Kraft und Richtig definieren, mit der der Feind weggestossen wird.

Weder Spieler noch Feind haben in Ninja Gaiden Sigma Frames in Animationen, wo sie keinen Schaden nehmen können, sogenante \textit{Invincibility-Frames}, ausser bei Spezialangriffen und Aufsteh-Animation. Denn das Spiel erlaubt weder den Spieler, noch die Feinde für eine lange Zeit ununterbrochen anzugreifen. Der Spieler ist von der Länge der Combos limitiert, wonach dieser eine Wiederholungs-Animation spielen muss, und die Feinde werden auch von Ihren Combos und den KI-Manager von übermässigen Angriffen limitiert. Das heisst, selbst wenn jemand hilfslos auf den Boden liegt und angegriffen werden kann, wird dies meistens nicht passieren. Der Charakter versucht, nach einer gewissen Zeiteinheit von Inaktivität automatisch aufzustehen.

\begin{description}
  \item[\textbf{M5}] \textbf{Wie hängen die gedrückten Tasten mit den visuellen Kampfereignissen zusammen, gibt es eine bestimmte Intention hinter den Tasten?}
\end{description}

Die Philosophie hinter den Angriffe in Ninja Gaiden folgt einen bestimmten Muster. Alle Aktionen beziehungsweise deren Animation, die auf den Knopf für den Leichten Angriff gebunden sind, haben einen schnelle Einführung in den Angriff, eine ebenfalls rasante Angriffsanimation und rapide Rückkehr zurück zum normalen Stand. Aktionen, die auf den Knopf für den Schweren Angriff gebunden sind, haben genau den Gegenteil als effekt. Deren Animationen sind oftmals länger und pompöuser, lassen den Protagonisten aber für eine längere Zeit verwundbar. Ausnahme zu diesen Regeln sind Angriffe, die am Ende von einer Combo ausgeführt werden, denn diese haben, wie erläutert, oftmals eine längere Wiederholungs-Animation. Ninja Gaiden schafft es aber, lange Angriffsanimationen nicht übermässig zu übertreiben. Genug um zu merken, dass man etwas schwierigeres gemacht hat, aber nicht genug, um den Spielfluss aufzuhalten.

Die erwähnte gefühlte Ungenauigkeit in der Ausführung der Combos liegt an einer speziellen Input-Buffer Implementierung, welches in Ninja Gaiden Sigma vorhanden ist, aber bei anderen ähnlichen Hack-and-Slash games der Era nicht in dieser Form existiert. Man kann in Ninja Gaiden bereits während der ersten Attacke alle Inputs der vollen Kombo in der richtigen Reihenfolge drücken. Diese werden dann automatisch ausgeführt. Hingegen besitzt \textit{Devil May Cry 3}, ein Hack-and-Slash Spiel aus dieser Era ein Input-Buffering System, welches den nächsten Angriff in einer Sequenz speichern kann. Dieser ist sogar jederzeit änderbar, solange die Aktion nicht bereits ausgeführt worden ist. Bei Ninja Gaiden kann das nicht, und weitere Tasteninputs erweitern entweder die Combo, falls diese noch erweitert werden kann, oder werden schlichtweg nicht registriert. Damit fühlt sich das System viel reaktionsschneller an, denn dies ist es in einen gewissen Sinne auch. Trotzdem kann man Argumentierung, dass die Implementierung von Input-Buffering von Ninja Gaiden Sigma aktiv gegen das Sinnlose drücken von Knöpfen agiert und damit für mehr geschicktes Spielen sorgt. Der Entwickler muss sich dann entscheiden, welche Route man gehen soll.


\end{document}
